%LaTeX-inställningar%%%%%%%%%%%%%%%%%%%%%%%%%%%%%%%%%%%%%%%%%%%%%%%%%%%
% Kompliera med pdflatex!!!
\documentclass[a4paper, 10pt]{article}

\usepackage{hyperref}
\usepackage[utf8]{inputenc}
\usepackage[T1]{fontenc}
\usepackage[swedish]{babel}
\usepackage{graphicx}
%\usepackage{ae} %För riktiga fonts?
\usepackage{fancyhdr}
\usepackage{lastpage}
\topmargin -20.0pt
\headheight 56.0pt
\setcounter{secnumdepth}{5}

%FYLL I VID ÄNDRINGAR%%%%%%%%%%%%%%%%%%%%%%%%%%%%%%%%%%%%%%%%%%%%%%%%%%
\newcommand{\updated}{2021-04-14} %Insert




%Dokumentstart%%%%%%%%%%%%%%%%%%%%%%%%%%%%%%%%%%%%%%%%%%%%%%%%%%%%%%%%%
\begin{document}
\pagestyle{fancy}

%Header%%%%%%%%%%%%%%%%%%%%%%%%%%%%%%%%%%%%%%%%%%%%%%%%%%%%%%%%%%%%%%%%
\renewcommand{\headrule}{\vbox to 0pt{\hfill\hbox to 292pt{\hrulefill}}}
\lhead{
\raisebox{-25pt}[0pt][10pt]{\includegraphics[width=60pt]{dteklogo.pdf}}
\parbox[b]{200pt}{
\textbf{Datateknologsektionen}\\
Chalmers studentkår\\
Ekonomiskt reglemente}}
\rhead{ \flushright
Sidan \thepage\ av \pageref{LastPage}\\
Uppdaterad \updated}

%Footer%%%%%%%%%%%%%%%%%%%%%%%%%%%%%%%%%%%%%%%%%%%%%%%%%%%%%%%%%%%%%%%%
\renewcommand{\footrulewidth}{\headrulewidth}
\lfoot{\flushleft Datateknologsektionen\\
    Rännvägen 8\\
    412 58 Göteborg}
\cfoot{}
\rfoot{ \flushright styret@dtek.se\\
www.dtek.se}
\newpage

%Titlepage%%%%%%%%%%%%%%%%%%%%%%%%%%%%%%%%%%%%%%%%%%%%%%%%%%%%%%%%%%%%%
\vspace*{\fill}
\begin{center}
{\Huge \textbf{Ekonomiskt reglemente för Datateknologsektionen}}\\
\includegraphics[width=300pt]{dteklogo.pdf}
\\{\LARGE Chalmers, Göteborg}
\end{center}
\vspace*{\fill}
\begin{center}
{\LARGE Uppdaterad: \updated}
\end{center}
\vspace*{\fill}


\newpage
\setcounter{tocdepth}{1}
\tableofcontents
\newpage

%Allmänt%%%%%%%%%%%%%%%%%%%%%%%%%%%%%%%%%%%%%%%%%%%%%%%%%%%%%%%%%%%%%%%
\section{Allmänt}
\subsection{}
Detta reglemente har skapats för att skapa tydlighet över vad sektionens pengar får och inte får användas till. Det är därmed tänkt att vara ett hjälpmedel för styrelsen, revisorerna och sektionens kommittéer i deras arbete. Grundfilosofin är att sektionens pengar är till för sektionens samtliga medlemmar, vilket i praktiken innebär att arrangemang som inte är öppna för dessa inte får belasta sektionens ekonomi. Viktigt att tänka på vid tolkandet utav reglementet är att det inte har utformats för att i onödan begränsa kommittéers möjligheter eller lust till roliga arrangemang och upptåg.
\subsection{Prisbasbelopp}
\begin{itemize}
  \item Det prisbasbelopp som skall användas under hela verksamhetsåret är det prisbasbelopp som är aktuellt vid verksamhetsårets början.
  \item Delar av prisbasbelopp skall avrundas till närmaste hundratal. 
\end{itemize}

%D-Styret%%%%%%%%%%%%%%%%%%%%%%%%%%%%%%%%%%%%%%%%%%%%%%%%%%%%%%%%%%%%%%
\section{D-Styret}
\subsection{}
D-Styret är ansvariga för sektionens hela ekonomi
\subsection{}
Ordförande är skyldig att
\begin{itemize}
  \item Teckna sektionens firma
  \item Se till att ordförande i varje kommitté har tillgång till och kunskap om sektionens stadgar, reglemente och ekonomiska reglemente 
\end{itemize}
\subsection{}
Kassören är skyldig att
\begin{itemize}
  \item Teckna sektionens firma
  \item Se till att kassör i varje kommitté har tillgång till och kunskap om sektionens stadgar, reglemente och ekonomiska reglemente 
  \item Fortlöpande kontrollera sektionens räkenskaper och bokföring
  \item I samråd med styrelsen upprätta budgetförslag till första ordinarie höstmötet
  \item Till varje sektionsmöte kunna redogöra för sektionens ekonomiska ställning
  \item Utbilda nya sektionsfunktionärer i hur sektionens bokförings och redovisningssystem skall användas
  \item Ansvara för styrelsens handkassa
\end{itemize}
\subsection{}
Styret har rättighet att inom ramen för styrets budget representera sektionen och ordna förtäring till styrets ledamöter i samband med styrelsemöten. Denna förtäring skall vara med som post i styrelsens budget.

%Sektionskommittéer%%%%%%%%%%%%%%%%%%%%%%%%%%%%%%%%%%%%%%%%%%%%%%%%%%%%
\section{Sektionskommittéer}
\label{sec:sektionsforeningar}
\subsection{}
Ordförande i sektionskommitté är skyldig att
\begin{itemize}
\item Kontinuerligt meddela kommitténs ekonomiska status till styret
\item Tillsammans med kommitténs kassör ansvara för att kommittén förvaltar sina tillgångar i enlighet med sektionens stadgar, reglementen och beslut.
\item Tillsammans med kommitténs kassör skriva ett bindande avtal med
styret angående skuldfrågan vid felaktig bokföring
\item Tillsammans med kommitténs kassör teckna kommitténs firma
\end{itemize}
\subsection{}
Kassör i sektionskommitté är skyldig att
\begin{itemize}
\item Föra kassabok av sådan typ som godkänts av sektionens revisorer
\item Tillsammans med kommitténs ordförande ansvara för att kommittén förvaltar sina tillgångar i enlighet med sektionens stadgar, reglementen och beslut.
\item Tillsammans med kommitténs ordförande skriva ett förbindande avtal med Styret angående skuldfrågan vid felaktigt förd bokföring
\item På varje sektionsmöte redovisa kommitténs ekonomiska situation
\item Arkivera sektionens bokföringar, på en plats anvisad av Styret, så lång tid som föreskrivs för den organisationsform som datateknologsektionen är.
\item Tillsammans med kommitténs ordförande teckna kommitténs firma
\end{itemize}
\subsection{}
Varje kommitté skall påbörja en ny kassabok i starten av sitt verksamhetsår
\subsection{Äskning}
Kommittéer som önskar extra ekonomiska medel till sin verksamhet skall inkomma med önskemål och skäl därför till styret. Styret kan bevilja extra medel om beloppet understiger 0,25 prisbasbelopp, om beloppet är större måste först sektionsmöte tillfrågas.
\subsection{Kapital}
\label{sec:sektionsforeningar_startkapital}
\subsubsection{Startkapital}
Kommittéer skall vid mandatperiodens början ha 0 prisbasbelopp i tillgångar. Kommittéer har rätt att låna 0,25 prisbasbelopp från sektionsstyrelsen i form avstartlån. DAG har rätt låna ytterligare 1 prisbasbelopp i form av utökat startlån.
\subsubsection{Slutkapital}
Samtliga kommitteers tillgångar och skulder tillfaller sektionsstyrelsen vid mandatperiodens slut.
\subsubsection{Stödkapital}
Styret har rätt att utföra kortfristiga lån till sektionskommittéer för att betala kristiska skulder i väntan på ankommande kapital.
\subsection{Budget}
För arrangemang med förväntad omsättning överstigande 0.75 prisbasbelopp skall en noggrann budget upprättas och godkännas av styret innan arrangemanget får genomföras.
\subsection{Arbetskläder}
Kommitté får lägga upp till 0.025 prisbasbelopp per person på arbetskläder. Rimliga arbetskläder kan anses vara en overall eller motsvarande, t-shirt samt annan
tröja/piké.
\subsection{Fri kost}
Vid arrangemang i kommitténs verksamhet får kommitténs ekonomi belastas med utgifter för inköp av fri kost, enligt följande tidsåtgång för arrangemanget, avrundat till närmaste femkrona:
\begin{itemize}
    \item 1–4 timmar: 0,00075 basbelopp per person.
    \item 4–8 timmar: 0,0015 basbelopp per person.
    \item Mer än 8 timmar: 0,00225 basbelopp per person.
\end{itemize}
Dessa utgifter får ej överskrida kommitténs ekonomiska kapacitet. Den fria kostens tillagning får inte medräknas i arrangemangets tidsåtgång

%Hobbykommittéer och andra grupper under sektionen%%%%%%%%%%%%%%%%%%
\section{Hobbykommittéer och andra grupper under sektionen}
\subsection{}
Ekonomiska medel för hobbykommittéer och andra grupper under sektionen skall regleras av styrets budget
\subsection{}
I de fall en hobbykommitté eller annan grupp under sektionen har beviljats egen ekonomi gäller även \S\ref{sec:sektionsforeningar}, med undantag för \S\ref{sec:sektionsforeningar_startkapital}, för dessa.

%Internrepresentation%%%%%%%%%%%%%%%%%%%%%%%%%%%%%%%%%%%%%%%%%%%%%%%%%%
\section{Internrepresentation}
\subsection{}
\label{sec:internreps}
Varje sektionskommitté och DNS har rättighet att varje verksamhetsår bekosta följande på sind ekonomi:
\begin{itemize}
    \item[-] En eller flera teambuildingaktiviteter för upp till 0,015 prisbasbelopp per person sammanlagt.
    \item[-] Överlämning för upp till 0,025 prisbasbelopp.
    \item[-] Aspning och aspplagg för upp till 0,065 prisbasbelopp sammanlagt.
\end{itemize}

Ovanstående får inte överskrida kommitténs ekonomiska resurser. Alkoholhaltiga drycker får ej bekostas med kommitténs pengar vid ovanstående arrangemang undantag gäller för teambuildingaktiviteter där alkoholhaltiga drycker får bekostastill upp till en tredjedel av teambuildingbudgeten.

\subsection{}
Styrelsen har rätt att bevilja hobbykommittéer eller andra grupperingar under teknologsektionen rättighet motsvarande \S\ref{sec:internreps} om det anses vara motiverat.

%Sponsring%%%%%%%%%%%%%%%%%%%%%%%%%%%%%%%%%%%%%%%%%%%%%%%%%%%%%%%%%%%%%
\section{Sponsring}
\subsection{}
Kommittéer får inte söka sponsring utan samråd med datas
arbetsmarknadsgrupp, DAG.

%Fonder%%%%%%%%%%%%%%%%%%%%%%%%%%%%%%%%%%%%%%%%%%%%%%%%%%%%%%%%%%%%%%%%
\section{Fonder}
De fonder som beskrivs i denna del syftar till sparande av medel för framtida användande enligt respektive fonds syfte, antingen planerat eller spontant. Ett uttag uren fond skall inte belasta Sektionsstyrelsens godkända budget, däremot skall inköp och investeringar vara föremål för redovisning i balans– respektive resultaträkning. Detta gäller om inget annat specificeras i fonden.
\subsection{Kapitalfonden}
\subsubsection{Syfte}
\label{sec:kapitalfond_syfte}
Syftet med kapitalfonden är att avlasta sektionens respektive de olika sektionskommittéernas ekonomi från stora investeringar. Pengarna skall användas till saker som har ett bestående värde, samt är till gagn för sektionens medlemmar direkt eller indirekt. Pengar skall ej användas till driftbidrag eller stöd för förgänglig verksamhet. Fonden skall ej användas till verksamhet som lokalfond är ämnad för.
\subsubsection{Förvaltning}
Fonden förvaltas av sektionens styrelse.
\subsubsection{Avsättning}
\begin{itemize}
\item En av styrelsen budgeterad summa som godkänts av sektionsmötet
\item All avkastning ifrån kapitalfonden tillförs kapitalfonden.
\end{itemize}
\subsubsection{Uttag}
\begin{itemize}
\item Sektionsstyrelsen har rätt att bevilja uttag ur fonden på belopp upp till totalt 0,25 prisbasbelopp per år. Uttag av belopp överstigande 0,25 prisbasbelopp skall godkännas av sektionsmöte innan medel utbetalas. Varje uttag skall redovisas inför nästkommande sektionsmöte.
\item Sektionsmötet äger rätt att vid välmotiverat behov göra uttag i strid med syftet i \S\ref{sec:kapitalfond_syfte}
\end{itemize}
\subsubsection{Revision}
Revidering sker varje år av sektionens revisorer. Revision sammanfaller med övrig på sektionen förekommande.

\subsection{Lokalfonden}
\subsubsection{Syfte}
\label{sec:lokalfond_syfte}
Syftet med lokalfonden är att säkra medel för underhåll och reparationer av sektionens lokaler, samt för inköp av inventarier. Pengarna skall användas till större reparationer och ommålningar av sektionslokalerna, samt möbler till trivselytor där teknologen i gemen har tillträde.
\subsubsection{Förvaltning}
Fonden förvaltas av sektionens styrelse.
\subsubsection{Avsättning}
\label{sec:lokalfond_avsattning}
\begin{itemize}
\item Tio (10) \% av under verksamhetsåret influtna sektionsavgifter tillförs lokalfonden.
\item En av styrelsen budgeterad summa som godkänts av sektionsmötet
\item All avkastning ifrån lokalfonden tillförs lokalfonden.
\end{itemize}
\subsubsection{Uttag}
\begin{itemize}
\item Sektionsstyrelsen disponerar 3/4 av årets tillförda kapital,
enligt första punkten i \S\ref{sec:lokalfond_avsattning}, för basdrift av sektionslokalerna. Sektionsstyrelsen har rätt att besluta om uttag på överstigande belopp, dock skall detta redovisas inför nästkommande sektionsmöte. Vid detta sektionsmöte skall i sådana fall även genomförda eller planerade inköp, reparationer och underhåll redovisas.
\item Sektionsmötet äger rätt att vid välmotiverat behov göra uttag i strid med syftet i \S\ref{sec:lokalfond_syfte}.
\end{itemize}
\subsubsection{Revision}
Revidering sker varje år av sektionens revisorer. Revision
sammanfaller med övrig på sektionen förekommande.



\subsection{Bilfonden}
\subsubsection{Syfte}
\label{sec:bilfond_syfte}
Syftet med bilfonden är att bygga upp en buffert för bilinköp.
\subsubsection{Förvaltning}
Fonden förvaltas av sektionens styrelse.
\subsubsection{Avsättning}
Två alternativ:
\begin{itemize}
\item 50 000 kronor per år.
\item En av styrelsens budgeterad summa som godkänts av sektionsmötet.
\item All avkastning ifrån bilfonden tillförs bilfonden.
\end{itemize}
\subsubsection{Uttag}
\begin{itemize}
\item Uttag får göras för inköp av en ny bil.
\item Uttag får också göras för omfattande reparation.
\item Sektionsmötet äger rätt att vid välmotiverat behov göra uttag i strid med syftet i \S\ref{sec:bilfond_syfte}.
\end{itemize}

\subsection{Idéfonden}
\subsubsection{Syfte}
\label{sec:idefond_syfte}
Syftet med idéfonden är att vara föremål för datateknologer som önskar medel tillarrangemang och tillförskaffningar vars effekt är till gagn för hela teknologsektionens medlemmar. Kommittéer eller andra grupper under teknologsektionen somönskarmedel för sådana ändamål är inte föremål för idéfonden om de har egenekonomi eller budget under sektionsstyrelsen.
\subsubsection{Förvaltning}
Idéfonden förvaltas av sektionsstyrelsen.
\subsubsection{Avsättning}
\begin{itemize}
\item En av styrelsen budgeterad summa som godkänts av sektionsmötet.
\item All avkastning ifrån idéfonden tillförs idéfonden.
\item Tio (10) \% av under verksamhetsåret influtna sektionsavgifter tillförs idéfonden.
\end{itemize}
\subsubsection{Uttag}
Sektionsstyrelsen har rätt att, på förslag av medlem i teknologsektionen, bevilja uttagpå upp till 0,25 prisbasbelopp som går i hand med idéfondens syfte. Sektionsmötet har rätt att, på förslag av medlem i teknologsektionen, bevilja uttagsom går i hand med \S\ref{sec:idefond_syfte}
\subsubsection{Motstridig uttag}
Sektionsmötet kan besluta om att göra uttag i strid med idéfondens syfte.

\end{document}
